\section{Method}
\label{sec:method}

\subsection{Problem Formulation}

Given a user $u$ with interaction history $H_u = \{h_1, h_2, \ldots, h_m\}$ and a candidate set $C = \{c_1, c_2, \ldots, c_n\}$ containing one positive item $c^+$ and $n-1$ negative items, the task is to rank candidates such that $c^+$ appears as high as possible.

In LLM-based recommendation, a language model $\mathcal{M}$ assigns a relevance score $s_i = f_\mathcal{M}(H_u, c_i)$ to each candidate. The ranking is determined by sorting candidates in descending order of $s_i$.

\subsection{Observation: LLM Confidence is Informative but Unreliable}

We first establish the motivating observation. When prompting Qwen3-8B to output a confidence score alongside its recommendation decision, we find:

\begin{enumerate}
    \item \textbf{Informative:} The confidence signal correlates with correctness (AUROC $> 0.5$ across all domains).
    \item \textbf{Unreliable:} The confidence is severely miscalibrated (ECE $\in [0.31, 0.66]$ for all baselines), with Maximum Calibration Error (MCE) exceeding 0.88 in every case.
    \item \textbf{Domain-dependent:} The same method's calibration quality varies unpredictably across domains (e.g., LLM2Rec AUROC ranges from 0.50 on Beauty to 0.70 on Books).
\end{enumerate}

This observation holds not only for our method but for all LLM-based recommendation baselines that produce numerical scores, establishing it as a \emph{backbone-level phenomenon} rather than a method-specific artifact.

\subsection{Task-Grounded Relevance Probability Scoring}

Instead of relying on implicit confidence or post-hoc calibration, we propose \textbf{task-grounded relevance probability scoring}: directly prompting the LLM to estimate the probability that a candidate is the user's next interaction, grounded in explicit reasoning about the user's history.

\paragraph{Prompt Design.} For each user-candidate pair $(u, c_i)$, we construct a structured prompt:

\begin{enumerate}
    \item Present the user's recent interaction history (up to 5 most recent items)
    \item Present the candidate item with title and description
    \item Ask the model to estimate $P(\text{relevant} | H_u, c_i) \in [0, 1]$
    \item Require a one-sentence justification (``reason'')
\end{enumerate}

The key design choice is asking for a \emph{task-specific probability} rather than a generic confidence score. The model must ground its estimate in observable evidence: category alignment, attribute matching, and purchase trajectory patterns.

\paragraph{Scoring Function.} The relevance probability output is used directly as the ranking score:
\begin{equation}
    s_i = P_\mathcal{M}(\text{relevant} | H_u, c_i)
\end{equation}

Candidates are ranked by $s_i$ in descending order. No post-hoc calibration, temperature scaling, or uncertainty penalty is applied.

\subsection{Why This Works}

The effectiveness of task-grounded relevance probability scoring stems from three properties:

\begin{enumerate}
    \item \textbf{Evidence grounding:} By requiring the model to justify its estimate, hallucinated high-confidence scores become less likely. The model must identify concrete evidence (category match, trajectory continuation) rather than pattern-matching surface features.

    \item \textbf{Probability framing:} Asking for a probability $\in [0,1]$ naturally produces a more calibrated output than asking for a ranking score or binary decision. The model's pretraining on probabilistic reasoning tasks transfers to this setting.

    \item \textbf{Task specificity:} The prompt explicitly frames the task as ``next interaction prediction'' rather than generic relevance. This aligns the model's output distribution with the evaluation metric (HR@$k$, NDCG@$k$).
\end{enumerate}

\subsection{Comparison with Existing Approaches}

Table~\ref{tab:method_comparison} contrasts our approach with existing LLM-based recommendation methods:

\begin{table}[h]
\centering
\small
\caption{Comparison of scoring approaches in LLM-based recommendation.}
\label{tab:method_comparison}
\begin{tabular}{lcccc}
\toprule
Method & Score Type & Calibrated & Evidence & Task-Specific \\
\midrule
LLMEmb & Embedding sim. & \xmark & \xmark & \xmark \\
LLM2Rec & Log-likelihood & \xmark & \xmark & Partial \\
IRLLRec & Intent score & \xmark & \xmark & Partial \\
ProEx/ProMax & Profile score & \xmark & \xmark & \cmark \\
RLMRec & Graph + LLM & \xmark & \xmark & \xmark \\
\midrule
\textbf{Ours} & Relevance prob. & \cmark & \cmark & \cmark \\
\bottomrule
\end{tabular}
\end{table}

Our method is the only approach that simultaneously produces calibrated scores, requires evidence grounding, and is task-specific. This combination explains its strong performance on domains where title/description information is sufficient for preference inference (Books, Electronics) and its competitive performance on domains with shorter, less informative item descriptions (Beauty, Movies).
